\begingroup
\renewcommand{\arraystretch}{1.45}
\clearpage\fancyhead[L]{\small REVIEW. DETERMINATION OF ATTRACTION}
\pdfbookmark[0]{REVIEW. DETERMINATION OF ATTRACTION}{work-native366}
\SourceLeaf{366}
\selectlanguage{english}
\begin{center}
\textbf{REVIEW.}
\end{center}
\begin{center}
\rule{0.62\linewidth}{0.4pt}\\[6pt]
Göttingische gelehrte Anzeigen.\quad 1818 February 9.\\[6pt]
\rule{0.62\linewidth}{0.4pt}
\end{center}

On January $17^{\mathrm{th}}$, Mr. Court Councillor \textsc{Gauss} presented a lecture to the Royal Society:

\begin{center}
\textit{Determinatio attractionis, quam in punctum quodlibet positionis datae exerceret
planeta, cuius massa per totam eius orbitam, ratione temporis, quo singulae partes
describuntur, uniformiter esset dispertita.}
\end{center}

\Needspace{27\baselineskip}
By a theorem which perhaps no one has yet expressly stated, but which is easily proved from the principles of physical astronomy, the secular variations of a planetary orbit caused by perturbation by another planet are the same whether the perturbing planet actually traverses its elliptical orbit according to Kepler's laws, or its mass is conceived as distributed over the circumference of the ellipse in such a way that portions of the ellipse otherwise traversed in equal times. receive equal shares of the whole mass: provided that the periods of revolution of the perturbed and perturbing planets are not in a rational ratio to one another. The problem forming the subject of this paper, namely the determination, which is not approximate but exact, not dependent on moderate eccentricity of the ellipse. but general, of the attraction exerted on any point in space by an elliptical ring of infinitesimally small and, according to the above
\SourceLeaf{367}
law, unchanging thickness, is therefore of great interest to physical astronomy. It is, however, no less remarkable from a purely mathematical point of view, because of the various devices required for its complete solution.

\Needspace{40\baselineskip}
It is not readily feasible to give an extract of the solution itself here. The author has chosen a purely analytical treatment; knowledgeable readers who follow it attentively will readily perceive the geometric counterparts of the individual quantities occurring in the investigation, and how it can be recast in geometric form. It may suffice here to give only the final result. Three unknown quantities $G,\ G',\ G''$ are determined by the roots of a cubic equation, whose nature permits it to be proved that it always has three real roots. The attraction of the elliptical ring resolved in any direction is then represented by an expression of the form $pP+qQ$, where $p$ and $q$ depend algebraically on $G,\ G',\ G''$, while $P$ and $Q$ denote the definite values of the integrals
\[
\int\frac{\cos T^{2}\,.\,\mathrm{d}T}
{2\pi(mm\cos T^{2}+nn\sin T^{2})^{\frac{3}{2}}}
\]
\[
\int\frac{\sin T^{2}\,.\,\mathrm{d}T}
{2\pi(mm\cos T^{2}+nn\sin T^{2})^{\frac{3}{2}}}
\]
when the integrations are taken from $T=0$ to $T=360^{0}=2\pi$, and where
\[
m=\surd(G+G'),\quad n=\surd(G+G'')
\]
\Needspace{45\baselineskip}
\par Since these integrals (except when $m=n$) are transcendental in nature. and are known to be connected with other much-studied transcendental quantities occurring in perturbation calculations, the solution could have been regarded as complete once it had been brought to this point. The author has, however, used this first opportunity that presented itself to give the first outlines of a \textit{new algorithm}, which he has used for many years to determine these transcendental quantities, and on which he will publish an extensive investigation in the future, leading to many remarkable results. Only the main theorems can be given here, omitting the proofs. If from two gi\SourceLeaf{368}%
ven positive quantities $m$ and $n$ other quantities $m',\ m'',\ m'''$ etc., $n',\ n'',\ n'''$ etc. are derived according to the following rules:
\[
\begin{array}{ll}
m'=\frac{1}{2}(m+n),&n'=\surd mn\\[0.4em]
m''=\frac{1}{2}(m'+n'),&n''=\surd m'n'\\[0.4em]
m'''=\frac{1}{2}(m''+n''),&n'''=\surd m''n''\ \mbox{etc.}
\end{array}
\]
\Needspace{27\baselineskip}
that is, if $m',\ n'$ are respectively the arithmetic and geometric means between $m$ and $n$; likewise $m'',\ n''$ the arithmetic and geometric means between $m'$ and $n'$, and so forth: then the terms both of the sequence $m,\ m',\ m'',\ m'''$ etc., and of the sequence $n,\ n',\ n'',\ n'''$ etc. approach a common limit $=\mu$ with extreme rapidity, which the author calls the arithmetic-geometric mean of $m$ and $n$. Evidently $\mu$ is at the same time the arithmetic-geometric mean of $m'$ and $n'$, or generally of any two corresponding terms of the two sequences. The author now proves that $\frac{1}{\mu}$ is the value of the integral
\[
\int\frac{\mathrm{d}T}{2\pi\surd(mm\cos T^{2}+nn\sin T^{2})}
\]
when the integration is taken from $T=0$ to $T=360^{0}$. It will readily be seen that this could also have been stated as follows: if the expansion of the function
\[
\frac{1}{\surd(\alpha+\beta\cos\psi)}
\]
gives the series
\[
A+B\cos\psi+C\cos2\psi+D\cos3\psi+\mathrm{etc.}
\]
then $\frac{1}{A}$ is always the arithmetic-geometric mean of the two quantities $\surd(\alpha+\beta)$ and $\surd(\alpha-\beta)$.

\Needspace{38\baselineskip}
A second equally important theorem is that, if the sum of the infinite series, which always converges very rapidly,
\[
2(m'm'-n'n')+4(m''m''-n''n'')+8(m'''m'''-n'''n''')+\mbox{etc.}
\]
in which the numerical coefficients form a geometric progression, is set $=\nobreak{}(mm-nn)\nu$, the value of the integral
\[
\int\frac{\cos2T\,.\,\mathrm{d}T}
{2\pi\surd(mm\cos T^{2}+nn\sin T^{2})}
\]
\SourceLeaf{369}
taken from $T=0$ to $T=360^{0}$ becomes $=-\frac{\nu}{\mu}$. Evidently the second coefficient of the above series is thereby also known, namely $B=-\frac{\nu}{2\mu}$, when $m=\surd(\alpha+\beta)$, $n=\surd(\alpha-\beta)$ have been set. All the following coefficients $C$, $D$, etc., however, are known to be determined algebraically and simply by the first two $A$ and $B$. A special, very convenient procedure for the numerical calculation of the quantities $m'm'-n'n'$, $m''m''-n''n''$ etc. is also taught in the paper itself.
\par
Application to the transcendental quantities of the present investigation finally gives the simple expressions
\[
P=\frac{1+\nu}{2mm\mu},\qquad Q=\frac{1-\nu}{2nn\mu}
\]
\Needspace{32\baselineskip}
\par Attentive readers will not fail to notice how many interesting problems connected with the transcendental quantities considered here are solved with the greatest ease by the algorithm explained. As an example, we give here the rectification of the ellipse. If its major semiaxis is set $=m$ and its minor semiaxis $=n$, its circumference becomes
\[
\begin{aligned}&=\frac{2\pi}{\mu}\bigl\{m'm'-2(m''m''-n''n'')\bigr.\\[5pt]&\qquad\bigl.-4(m'''m'''-n'''n''')-8(m''''m''''-n''''n'''')-\mathrm{etc.}\bigr\}\end{aligned}
\]
Another example is the duration of pendulum oscillations over finite arcs, which is to the duration of infinitesimally small oscillations as unity is to the arithmetic-geometric mean between $1$ and the cosine of one quarter of the whole arc of oscillation.
\Needspace{14\baselineskip}
\par
Finally, it must be mentioned that the author believed it necessary to present these results in their original form, as he had found them many years ago independently of similar investigations by \textsc{Lagrange} and \textsc{Legendre}, although they could in part easily have been derived from the discoveries of these geometers; partly because that form seemed to him to have substantial advantages, and partly because in just this form they constitute the beginning of a much more extensive theory, in which his work has taken a direction quite different from that of the geometers named.
\begin{center}\rule{0.25\linewidth}{0.4pt}\end{center}
\clearpage
\endgroup
